% Generated by tools/build_v6.py; edit the assembly rules, not this file.
\section{An exact dual for the robust stopping program}
\label{sec:exact-dual}

Triangular stopping representations and their adaptive comparison have
close antecedents in Lee--Valiant \cite{leevaliant}. Here the sign allocation
ranges over the full convex rarity class, and bias, variance, and work
are robust constraints. We derive the exact four-prior dual of this
program and absorb its variance and work references into the endpoints.
Appendix~\ref{sec:flow-foundations} proves the stopping-flow representation
and its observable realization in full.

Throughout this section, $\mu$ is any probability measure on $(0,1/2]$,
$0<\rho\le1/4$, $B>0$, and $N\ge1$ is an integer. Put
\[
 \mathcal Q_\rho=\{q\in L^\infty(\mu):0\le q\le1,\
                \theta(q)=\textstyle\int q\,d\mu\le\rho\}.
\]
For a state $s=(n,k)$, $0\le k\le n\le N$, define
\begin{align}
 m_{nk}(q)&=\int\binom nk\bigl[
 (1-q(d))(1/2-d)^k(1/2+d)^{n-k}\notag\\[-2pt]
 &\hspace{39mm}+q(d)(1/2+d)^k(1/2-d)^{n-k}\bigr]\,\mu(dd).
 \label{eq:dual-count-coefficients}
\end{align}
In particular $m_{00}(q)=1$. These are probabilities of count states at
fixed depth, before imposing a stopping rule.

\subsection{The primal and its four-prior dual}

For clarity, the finite program being dualized is recalled in compact
form. Its nonnegative stopping and continuation arrays $a,c$, real output
array $z$, and nonnegative moment array $t$ obey
\begin{align}
 &a_{00}=z_{00}=t_{00}=0,\qquad c_{00}=\lambda\ge0,
       \qquad c_{Nk}=0,\notag\\
 &a_{nk}+c_{nk}=\frac{n-k}{n}c_{n-1,k}
                    +\frac{k}{n}c_{n-1,k-1},
       \qquad 1\le n\le N,\notag\\
 &z_{nk}^2\le a_{nk}t_{nk}.
 \label{eq:dual-flow-cone}
\end{align}
Missing indices are zero. These are normalized prefix-weight flows, not
unconditional state probabilities. Write $\mathcal C_N$ for this closed
convex cone. The squared program value is
\begin{equation}
\begin{split}
 P_N(B)^2=\min\;&b^2+v\quad
       \text{over }(a,c,z,t,\lambda)\in\mathcal C_N,\ c_0\in\R,\ b,v\ge0,\\
 \text{subject to }&
 \sup_{q\in\mathcal Q_\rho}
       \bigl|c_0+\langle m(q),z\rangle-\theta(q)\bigr|\le b,\\
 &\sup_q\langle m(q),t\rangle\le v,\qquad
   \sup_q\langle m(q),c\rangle\le B.
\end{split}
\label{eq:dual-primal}
\end{equation}
Stopping sums exclude the root and continuation sums exclude depth $N$.
Every feasible point realizes a Poisson estimator with uniform RMS at
most $\sqrt{b^2+v}$ and expected work at most $B$
(Proposition~\ref{prop:ff-realization}).

Use the closed perspective
\[
 \mathfrak r(x,y)=
 \begin{cases}
 x^2/y,&y>0,\\
 0,&x=y=0,\\
 +\infty,&y=0, x\ne0.
 \end{cases}
\]
Given four priors $q_L,q_H,q_v,q_w\in\mathcal Q_\rho$, put
$d_s=m_s(q_H)-m_s(q_L)$, $v_s=m_s(q_v)$, $w_s=m_s(q_w)$, and define
\begin{equation}
 K_N(q_L,q_H;q_v,q_w)
 =\sup_{A:\,1\le T_A\le N}
 \frac{\displaystyle\sum_s a_s^A\mathfrak r(d_s,v_s)}
      {\displaystyle\sum_s c_s^A w_s}.
 \label{eq:dual-information-ratio}
\end{equation}
Here $A$ is a one-unit policy that always requests its first observation;
its normalized flow has root intensity one. A zero stopping coefficient
contributes zero even when the perspective is infinite. The denominator
is $\E_{q_w}T_A\in[1,N]$.

For a signed measure $P_H-P_L$, write
\[
 \mathcal I(P_H,P_L;P_v)
 =\int\left(\frac{d(P_H-P_L)}{dP_v}\right)^2\,dP_v,
\]
with value infinity unless the signed measure is absolutely continuous.
The numerator in \eqref{eq:dual-information-ratio} equals
$\mathcal I(P_H^A,P_L^A;P_v^A)$. Include the independent seed in the
transcript: its likelihood ratio depends only on the terminal count
state, so this identity covers order-dependent stopping and randomized
policies as well.

\begin{theorem}[Exact four-prior dual]\label{thm:exact-dual}
The minimum in \eqref{eq:dual-primal} is attained, and
\begin{equation}
 \boxed{
 P_N(B)^2=
 \max_{q_L,q_H,q_v,q_w\in\mathcal Q_\rho}
 \frac{[\theta(q_H)-\theta(q_L)]_+^2}
      {4+B K_N(q_L,q_H;q_v,q_w)}.}
 \label{eq:exact-four-prior-dual}
\end{equation}
The ratio is zero when $K_N=\infty$. The maximum is attained with positive
value, positive target separation, and finite $K_N$. Swapping the endpoints
allows the positive part to be replaced by an absolute value.
\end{theorem}

\begin{proof}
\emph{Compactness and constraint qualification.}
The finite moment image
$\{(\theta(q),m(q)):q\in\mathcal Q_\rho\}$ is compact and convex:
$\mathcal Q_\rho$ is weak-star compact, and every coordinate is a bounded
linear functional of $q$. All robust functions in the program are thus
finite continuous convex functions.

The zero-flow output $\rho/2$ gives objective $\rho^2/4$. On this objective
sublevel, $b\le\rho/2$, $v\le\rho^2/4$, and the root work charge gives
$\lambda\le B$. The normalized recursion implies $0\le a_s,c_s\le B$.
For the constant prior $q_*=\rho/2$, $m_s(q_*)>0$ unless state $s$ is
impossible under every prior. At every nonnull stopping state,
\[
 t_s\le\frac{\rho^2}{4m_s(q_*)},\qquad |z_s|\le\sqrt{Bt_s}.
\]
At universally null states set $z_s=t_s=0$, without changing any mean,
moment, or cost. The bias constraint at $q=0$ then bounds $c_0$.
Consequently this feasible sublevel, with null output coordinates removed,
is compact. The primal minimum is attained.

There is also a strict relative-interior feasible point. Choose
$0<\lambda<B/4$, continue with probability $1/2$ at each nonterminal
observed state, and stop at depth $N$. All free stopping and continuation
flows are positive and expected work is less than $2\lambda<B$.
Set $z=0$, take all free $t_s>0$, and choose sufficiently large $b,v$.
Every perspective inequality and all four robust inequalities are strict.
Finite-dimensional convex Slater duality therefore gives no duality gap,
attained inequality multipliers, and KKT conditions at a primal optimum.

\emph{Reduction to four priors.}
Split the absolute bias constraint into its two one-sided inequalities.
Each robust function is a maximum of affine functions over a compact
convex set. A subgradient is a convex combination of active affine
gradients. Since $q\mapsto(\theta(q),m(q))$ is affine and
$\mathcal Q_\rho$ is convex, this combination is represented by one active
prior. Thus the two bias subgradients, the variance subgradient, and the
work subgradient may be represented respectively by
$q_L,q_H,q_v,q_w$.

Let their nonnegative multipliers be $\alpha_L,\alpha_H,\beta,\eta$.
The free offset forces $\alpha_L=\alpha_H=\alpha$. Minimizing
$b^2-2\alpha b$ over $b\ge0$ gives $-\alpha^2$, and minimizing
$(1-\beta)v$ over $v\ge0$ requires $\beta\le1$. The remaining expression
on the flow cone is
\begin{equation}
 \beta\langle m(q_v),t\rangle
 -\alpha\langle m(q_H)-m(q_L),z\rangle
 +\eta\langle m(q_w),c\rangle.
 \label{eq:dual-cone-functional}
\end{equation}
Homogeneity requires this expression to be nonnegative everywhere on
$\mathcal C_N$; otherwise scaling a negative point gives infimum minus
infinity. When it is nonnegative its infimum is zero. The dual objective is
\begin{equation}
 \alpha[\theta(q_H)-\theta(q_L)]-\alpha^2-\eta B.
 \label{eq:dual-price-objective}
\end{equation}

We may take $\beta=1$. Indeed, for any $q_0\in\mathcal Q_\rho$, replace
$\beta m(q_v)$ by
\[
 m\bigl(\beta q_v+(1-\beta)q_0\bigr)
 =\beta m(q_v)+(1-\beta)m(q_0).
\]
The additional coefficient on $t$ is nonnegative, so cone feasibility is
preserved while the dual objective is unchanged. This also preserves
attainment.

\emph{Eliminating outputs and the work price.}
For $\alpha>0$, the perspective constraint gives, at fixed stopping flow,
\[
 \inf_{z_s^2\le a_st_s}
       \{v_st_s-\alpha d_sz_s\}
 =-\frac{\alpha^2a_s}{4}\mathfrak r(d_s,v_s).
\]
For positive $v_s$ this is completion of the square. If $a_s=0$ the
contribution is zero; if $a_s>0$, $v_s=0$, and $d_s\ne0$, the infimum is
minus infinity. The remaining zero case has value zero.
Every nonzero stopping/continuation flow scales to root intensity one.
At $\lambda=0$, $a=c=z=0$ and surplus $t$ has a nonnegative contribution.
Thus nonnegativity of
\eqref{eq:dual-cone-functional} is exactly
\[
 \eta\ge\alpha^2 K_N/4.
\]
For finite $K_N$, optimizing \eqref{eq:dual-price-objective} yields
\[
 \sup_{\alpha\ge0}
 \left\{\alpha\Delta\theta
    -\alpha^2(1+BK_N/4)\right\}
 =\frac{(\Delta\theta)_+^2}{4+BK_N}.
\]
When $K_N=\infty$, only the zero-value choice $\alpha=\eta=0$ remains.
Weak duality applies to every quadruple, and the attained KKT dual above
supplies an optimal quadruple. This proves the identity.

Finally take $q_L=0$ and $q_H=q_v=\rho$ constant. The divergence of the
latent sign contrast relative to $q_v$ is $\rho/(1-\rho)$. Data processing
bounds every stopped-transcript numerator by this quantity, and work is
at least one. Thus, for any $q_w$,
\begin{equation}
 P_N(B)^2\ge\frac{\rho^2}{4+B\rho/(1-\rho)}>0.
 \label{eq:dual-positive-value}
\end{equation}
An attained optimal dual therefore has $\alpha>0$, finite $K_N$, and
positive target separation, as asserted.
\end{proof}

\subsection{A Bellman recursion with explicit null-state conventions}

Divide each coefficient in \eqref{eq:dual-count-coefficients} by
$\binom nk$ and write $p_q(n,k)$ for the probability of one particular
ordered prefix with that count. These coefficients satisfy
\[
 p_q(n,k)=p_q(n+1,k)+p_q(n+1,k+1).
\]
In this subsection let $d_s=p_H(s)-p_L(s)$, $v_s=p_v(s)$,
$w_s=p_w(s)$, and $r_s=\mathfrak r(d_s,v_s)$ denote the corresponding
\emph{prefix} coefficients. At the root $v_0=w_0=1$ and $d_0=r_0=0$.

\begin{proposition}[Bellman representation]\label{prop:dual-bellman}
Suppose $K_N<\infty$. For $\kappa\ge0$, set
\begin{align}
 U_{Nk}(\kappa)&=r_{Nk},\notag\\
 U_{nk}(\kappa)&=\max\{r_{nk},
 U_{n+1,k}(\kappa)+U_{n+1,k+1}(\kappa)-\kappa w_{nk}\},
 \quad 1\le n<N.
 \label{eq:dual-bellman}
\end{align}
Then
\begin{equation}
 \Phi_N(\kappa):=U_{10}(\kappa)+U_{11}(\kappa)-\kappa
 =\max_{1\le T_A\le N}
 \{\mathcal I(P_H^A,P_L^A;P_v^A)-\kappa\E_{q_w}T_A\}.
 \label{eq:dual-price-value}
\end{equation}
It is finite, continuous, piecewise affine, and strictly decreasing, and
its unique nonnegative zero is $K_N$. A deterministic count-state policy
attains the maximum at every price. If some prefix has $v_s=0$ and
$d_s\ne0$, then $K_N=\infty$ instead.
\end{proposition}

\begin{proof}
At any observed prefix, stopping earns $r_s$ and continuation earns the
two child values minus the work charge $\kappa w_s$. Backward induction
therefore proves \eqref{eq:dual-bellman}. Root activation is compulsory,
which gives the final charge $\kappa$. All coefficients depend only on
the count, so the maximizing choice may be deterministic and depend only
on $(n,k)$. Randomized or order-dependent choices cannot improve it.
There are finitely many such deterministic policies, each with work in
$[1,N]$. Their affine price values have slopes in $[-N,-1]$, proving the
continuity and strict monotonicity claims. The zero-price value is
nonnegative, the value tends to minus infinity, and the zero is precisely
the maximal information-to-work ratio. If a prefix has zero reference
likelihood and nonzero contrast, the policy stopping at its depth has
infinite information.
\end{proof}

The recursion also admits an information-increment form. Put
\begin{equation}
 j_{nk}=r_{n+1,k}+r_{n+1,k+1}-r_{nk}\ge0.
 \label{eq:dual-information-increment}
\end{equation}
Nonnegativity is Cauchy--Schwarz. At a zero-reference state, finite $K_N$
forces both its contrast and its children's contrasts to vanish, so
$j_s=0$ there. If $b_s^A$ is the unnormalized continuation-prefix flow,
telescoping the flow identities gives
\begin{equation}
 \mathcal I(P_H^A,P_L^A;P_v^A)=\sum_{n<N,k}b_s^A j_s,
 \qquad \E_{q_w}T_A=\sum_{n<N,k}b_s^A w_s.
 \label{eq:dual-increment-telescope}
\end{equation}
Subtracting $r_s$ from the Bellman values gives terminal values $V_N=0$
and
\[
 V_s(\kappa)=\max\{0,j_s-\kappa w_s+V_{s0}(\kappa)+V_{s1}(\kappa)\},
 \qquad
 \Phi_N(\kappa)=j_0-\kappa+V_{10}(\kappa)+V_{11}(\kappa).
\]
Thus the ratio accounts for both information gained at a state and the
work needed to reach it. The pointwise inequalities $j_s\le\kappa w_s$
are sufficient, though generally not necessary, for $K_N\le\kappa$.

When $q_v$ dominates $q_w$ on all relevant prefixes, the posterior form is
\[
 \Phi_N(\kappa)=\sup_{1\le T\le N}
 \E_{q_v}\left[L_T^2-\kappa\sum_{n<T}W_n\right],
 \qquad L_n=d_n/v_n,\quad W_n=w_n/v_n.
\]
The process $L_n$ is a signed-likelihood martingale, initially zero. The
unnormalized Bellman recursion itself requires no such domination
assumption, and is the appropriate form at degenerate states.

\section{A midpoint modulus and a sharper adaptive comparison}
\label{sec:midpoint-dual}

Define the two-prior quantity
\begin{equation}
 D_N(B)^2=
 \sup_{q_-,q_+\in\mathcal Q_\rho}
 \frac{|\theta(q_+)-\theta(q_-)|^2}
 {4+B K_N(q_-,q_+;\bar q,\bar q)},
 \qquad \bar q=(q_-+q_+)/2.
 \label{eq:dual-midpoint-modulus}
\end{equation}

\begin{theorem}[Absorbing the two references]\label{thm:midpoint-modulus}
For every law, budget, and finite cap in the stated range,
\begin{equation}
 D_N(B)\le P_N(B)\le3D_N(B).
 \label{eq:dual-midpoint-comparison}
\end{equation}
\end{theorem}

\begin{proof}
The first inequality restricts the four-prior dual. For the second, start
with any quadruple of finite $K_N$ and set
\[
 q'_-=(q_L+q_v+q_w)/3,\qquad q'_+=(q_H+q_v+q_w)/3.
\]
Their midpoint dominates $q_v$ and $q_w$ in the following precise
experiment sense: for every policy,
$P_{\bar q'}^A\ge P_{q_v}^A/3$ and
$\E_{\bar q'}T_A\ge\E_{q_w}T_A/3$.
The signed transcript contrast is the old contrast divided by three.
Consequently its information numerator is at most one third of the old
numerator, and
\[
 K_N(q'_-,q'_+;\bar q',\bar q')
 \le K_N(q_L,q_H;q_v,q_w).
\]
The target gap is divided by three. Substitution in
\eqref{eq:dual-midpoint-modulus} and then
\eqref{eq:exact-four-prior-dual} proves the result.
\end{proof}

For midpoint references, the numerator is twice triangular
discrimination. With $h(P,Q)=1-\operatorname{Aff}(P,Q)$,
\begin{equation}
 4h(P_+,P_-)
 \le 2\int\frac{(p_+-p_-)^2}{p_++p_-}
 =\mathcal I(P_+,P_-;(P_++P_-)/2)
 \le8h(P_+,P_-).
 \label{eq:dual-hellinger-comparison}
\end{equation}
These inequalities follow by factoring
$(p_+-p_-)^2=(\sqrt{p_+}-\sqrt{p_-})^2
(\sqrt{p_+}+\sqrt{p_-})^2$ and bounding the second factor divided by
$p_++p_-$ between one and two. Hence the midpoint modulus is also a
Hellinger-information-per-work modulus, up to universal constants. It
still involves an optimal local stopping rule; it is not a closed
tail-moment or entropy formula.

\begin{lemma}[Squared-loss testing by affinity]\label{lem:affinity-testing}
If two terminal laws have target gap $\Delta$, every estimator has worst
RMS at least $\Delta\operatorname{Aff}(P_-,P_+)/2$.
\end{lemma}

\begin{proof}
The maximum MSE is at least the average MSE for equally likely endpoints.
The minimum average MSE is
\[
 \frac{\Delta^2}{2}\int\frac{p_-p_+}{p_-+p_+}.
\]
Cauchy--Schwarz gives
$\operatorname{Aff}(P_-,P_+)^2
\le2\int p_-p_+/(p_-+p_+)$, proving the claim.
\end{proof}

Let $R_N^{\mathrm{tot}}(B)$ denote population minimax RMS with at most
$N$ total observations and uniform expected work at most $B$. This total
cap is different from the per-unit cap used in $P_N$.

\begin{theorem}[A direct adaptive comparison]\label{thm:dual-adaptive}
\begin{equation}
 R_N^{\mathrm{tot}}(B)\ge\frac38D_N(B),\qquad
 P_N(B)\le8R_N^{\mathrm{tot}}(B).
 \label{eq:dual-capped-adaptive}
\end{equation}
\end{theorem}

\begin{proof}
Fix a pair in \eqref{eq:dual-midpoint-modulus}, with target gap $\Delta$
and information ratio $K$. Interpolate its endpoints toward their
midpoint by $s\in[0,1]$. The midpoint stays fixed, the target gap becomes
$s\Delta$, and the information ratio becomes exactly $s^2K$.

For every induced local policy of a cap-$N$ population algorithm,
\eqref{eq:dual-hellinger-comparison} gives
\[
 h(P_-^A,P_+^A)\le\frac{s^2K}{4}\E_{\bar q}T_A.
\]
This inequality also covers policies that sometimes take zero
observations. Their zero-work event depends only on the seed and has zero
signed contrast; conditioning on activation gives a root-active policy,
and both information and expected work scale by the activation
probability. Suppressing the seed can only decrease information.

Apply Lemma~\ref{lem:tf-midpoint}.
It realizes each hybrid as a one-genuine-unit policy, simulating other
units from independent midpoint tapes in its auxiliary seed. The hybrid
global transcript is a function of that local transcript, so data
processing applies even if the seed reveals the other tapes. Summing the
midpoint local work counts every population observation exactly once.
Therefore
\[
 h_{\mathrm{pop}}\le s^2BK/4.
\]
If $K=0$ take $s=1$; otherwise take $s=\min\{1,(BK)^{-1/2}\}$.
Then $h_{\mathrm{pop}}\le1/4$. Lemma~\ref{lem:affinity-testing} yields
\[
 R_N^{\mathrm{tot}}(B)\ge\frac38s\Delta
 \ge\frac38\frac\Delta{\sqrt{4+BK}}.
\]
Taking the pair supremum proves the first claim; the second follows from
Theorem~\ref{thm:midpoint-modulus}.
\end{proof}

\begin{corollary}[An improved finite characterization]
\label{cor:dual-finite-characterization}
For every integer $N\ge1$,
\begin{equation}
 \left(\frac{P_N(B)}8-
       \frac\rho2\sqrt{\frac B{N+1}}\right)_+
 \le R_{\mu,\rho}(B)\le P_N(B).
 \label{eq:dual-finite-certificate}
\end{equation}
At the explicit cap
\[
 N_*=\left\lceil64\max\{B,\rho B^2\}\right\rceil,
\]
this gives
\begin{equation}
 \boxed{R_{\mu,\rho}(B)\le P_{\mu,\rho,N_*}(B)
                         \le16R_{\mu,\rho}(B).}
 \label{eq:dual-factor-sixteen}
\end{equation}
\end{corollary}

\begin{proof}
Clip any estimator to
$[0,\rho]$ and, before requesting observation $N+1$, returns $\rho/2$.
Markov's inequality and the $L^2$ triangle inequality give
\[
 R_N^{\mathrm{tot}}(B)
 \le R_{\mu,\rho}(B)+\frac\rho2\sqrt{\frac B{N+1}}.
\]
Combining this with Theorem~\ref{thm:dual-adaptive} and the constructive
flow upper bound proves \eqref{eq:dual-finite-certificate}.
Lemma~\ref{lem:tf-floor} gives
\[
 R_{\mu,\rho}(B)\ge\frac1{16}
                 \min\{\rho,\sqrt{\rho/B}\}.
\]
At $N=N_*$ the truncation remainder is at most that floor. Thus
$P_{N_*}(B)\le8(R+R)=16R$.
\end{proof}

The program remains a uniform risk certificate rather than an exact formula
for minimax risk, and its finite dimension does not imply an efficient
algorithm at every budget. Theorem~\ref{thm:entropy-tail-law} supplies a
complementary additive expression involving only the magnitude law, with
its matching lower proof in Appendix~\ref{sec:matching-lower}. Its
conservative scalar constants are distinct from this program factor.


\begin{remark}[Scope of the factor-sixteen guarantee]
The comparison uses the sufficient cap in
Corollary~\ref{cor:dual-finite-characterization}, or a cap that passes
the verified genie stopping condition in Section~\ref{sec:genie}.
It is not a comparison for every fixed cap. In particular, an ordinary
fixed-cap dual alone is not a lower certificate for unrestricted risk.
With verified bounds $L\le G_N$ and $P_N\le U$, the orientation is
\[
 L\le G_N\le P_\infty\le P_N\le U,\qquad L/8\le R\le U.
\]
Thus $U\le2L$ implies $U\le16R$.
\end{remark}
